% \usepackage[bookmarks,colorlinks,breaklinks]{hyperref}
\hypersetup{urlcolor=blue, colorlinks=true, citecolor=green!50!black, linkcolor=blue}
\usepackage{multirow}
\usepackage[T1]{fontenc} % For polish hook https://tex.stackexchange.com/questions/75396/how-to-use-polishhook-symbol
\usepackage{amsmath}
\usepackage{amsfonts}
\usepackage{amssymb}
\usepackage{amsthm}
\usepackage{thmtools}
\usepackage{pgffor}
\usepackage{xcolor}
% \usepackage[lined,boxed,ruled,norelsize,algo2e,linesnumbered,noend]{algorithm2e}
\usepackage{cleveref}
\usepackage{graphicx}
\usepackage{geometry}
% \geometry{verbose,tmargin=1in,bmargin=1in,lmargin=1in,rmargin=1in}
\usepackage{enumitem}


% \newtheoremstyle{mystyle}
%   {12pt} % Space above
%   {24pt} % Space below
%   {} % Body font
%   {} % Indent amount
%   {\bfseries} % Theorem head font
%   {.} % Punctuation after theorem head
%   {.5em} % Space after theorem head
%   {} % Theorem head spec (can be left empty, meaning `normal')

% \theoremstyle{mystyle}


\newtheorem{assumption}{Assumption}
\newtheorem{definition}{Definition}
\newtheorem{theorem}{Theorem}
\newtheorem{lemma}{Lemma}
\newtheorem{proposition}{Proposition}
\newtheorem{corollary}{Corollary}

\usepackage{thmtools}

% \declaretheorem[refname={Theorem,Theorems},style=mystyle, Refname={Theorem,Theorems}]{theorem}
% \declaretheorem[numberlike=theorem, style=mystyle]{lemma}
% \declaretheorem[numberlike=theorem, style=mystyle]{invariant}
% \declaretheorem[numberlike=theorem, style=mystyle]{corollary}
% \declaretheorem[numberlike=theorem, style=mystyle]{definition}
% \declaretheorem[numberlike=theorem, style=mystyle]{claim}
% \declaretheorem[numberlike=theorem, style=mystyle]{fact}
% \declaretheorem[numberlike=theorem, style=mystyle]{assumption}
% \declaretheorem[numbered=no]{remark}

% number sets
\newcommand{\Real}{\mathbb{R}}
\newcommand{\Nat}{\mathbb{N}}
\newcommand{\F}{\mathbb{F}}
\newcommand{\Z}{\mathbb{Z}}
\renewcommand{\S}{\mathbb{S}}%

% probability
\renewcommand{\P}{\mathbb{P}}
\newcommand{\E}{\mathbb{E}}

% redfine some classical notation
\renewcommand{\tilde}{\widetilde}
\renewcommand{\hat}{\widehat}
\renewcommand{\bar}{\overline}

% operators
\newcommand{\median}{\operatorname{median}}

\DeclareMathOperator*{\argmin}{argmin}
\DeclareMathOperator*{\argmax}{argmax}
\newcommand{\poly}{\operatorname{poly}}
\newcommand{\polylog}{\operatorname{polylog}}
\newcommand{\mdiag}{\mathbf{Diag}} % diagonal matrix with vector on diag
\newcommand{\sign}{\operatorname{sign}}%
\newcommand{\nnz}{\operatorname{nnz}}%
\newcommand{\tail}{\operatorname{tail}}%
\newcommand{\dist}{\operatorname{dist}}%

% matrices
% \mX is \mathbf{X}
\foreach \x in {A,...,Z}{%
	\expandafter\xdef\csname m\x\endcsname{\noexpand\mathbf{\x}}
}
% \omX is \overline{\mathbf{X}}
\foreach \x in {A,...,Z}{%
	\expandafter\xdef\csname om\x\endcsname{\noexpand\overline{\noexpand\mathbf{\x}}}
}

\foreach \x in {A,...,Z}{%
	\expandafter\xdef\csname c\x\endcsname{\noexpand\mathcal{\x}}
}




\newcommand{\ob}{\overline{b}}
\newcommand{\ox}{\overline{x}}
\newcommand{\os}{\overline{s}}


\usepackage{mathtools}


% Reduce paragraph spacing https://tex.stackexchange.com/questions/4891/how-do-i-control-the-spacing-above-a-new-paragraph
\makeatletter
\renewcommand{\paragraph}{%
	\@startsection{paragraph}{4}%
	{\z@}{1.25ex \@plus 1ex \@minus .2ex}{-1em}%
	{\normalfont\normalsize\bfseries}%
}
\makeatother

\def\norm#1{\left\| #1 \right\|}



% \usepackage[
%     sorting=nyt,
% 	style=alphabetic,
%     maxalphanames=3,
%     minalphanames=3,
%     maxcitenames=3,
%     minnames=1,
%     maxbibnames=20,
% 	backref=true,
% 	doi=true,
% 	url=true,
% 	backend=biber
% ]{biblatex}
\usepackage{tcolorbox}


\newcommand{\thmbox}[1]{%
  \begin{tcolorbox}[
      colback = blue!5!white,
      colframe = blue!75!black,
      width          = \dimexpr\textwidth + 1cm\relax,
      % width    = \textwidth,   % full text width
      % left     = 5pt,          % no extra padding on the left
      % right    = 5pt,           % no extra padding on the right
      enlarge left by  = -0.5cm,
      enlarge right by = -0.5cm
  ]
    #1
  \end{tcolorbox}%
}

\newcommand{\mainthmbox}[1]{%
  \begin{tcolorbox}[colback=blue!5!white,colframe=blue!75!black,boxrule=3pt,
    ]
    #1
  \end{tcolorbox}%
}


\newcommand{\defbox}[1]{%
  \begin{tcolorbox}[colback=orange!5!white,colframe=orange!75!black, 
    width          = \dimexpr\textwidth + 1cm\relax,
    % width    = \textwidth,   % full text width
    % left     = 5pt,          % no extra padding on the left
    % right    = 5pt,           % no extra padding on the right
    enlarge left by  = -0.5cm,
    enlarge right by = -0.5cm]
    #1
  \end{tcolorbox}%
}

\newcommand{\blockbox}[1]{%
  \begin{tcolorbox}[colback=yellow!5!white,colframe=yellow!75!black]
    #1
  \end{tcolorbox}%
}

\newcommand{\shortminus}{\scalebox{0.5}[1.0]{\( - \)}}

% \usepackage{mathabx}  

\usepackage{cancel}
% \usepackage{arydshln}

\usepackage{geometry}
 \geometry{
        a4paper,
        total={170mm,257mm},
        left=20mm,
        right=10mm,
        top=20mm,
        }
        
% \usepackage{graphicx} % Required for inserting images

\usepackage{tikz}

\usetikzlibrary{arrows.meta,
                positioning,
                shadows}
                
\usetikzlibrary{shapes,snakes}
\usepackage{bm}
\usepackage{xspace}

\usepackage{cancel}


\setlist[itemize]{topsep=1pt, itemsep=1pt, parsep=1pt, leftmargin=*}

\usepackage{subcaption}
% \usepackage{graphicx}

